\section{Open and Responsible：权重开放只是四分之一}

这一章把 data governance 放在模型 release 的核心。开放性至少包含可访问数据、可访问模型、可复现研究和完整文档；responsibility 还要求 opt-out、PII、license、evaluation 与 limitation reporting。

\subsection{四层开放：Data、Models、Research、Documentation}

slide 将开放拆成四列。数据需要 inspection/opt-out/PII；模型需要 weights、training code、evaluation interface；研究需要 processing/training/evaluation reproducibility；文档需要 dataset/model cards、governance、architecture、cost 与 limitations。

\lecturefigure{slide-61.jpg}{Open and Responsible Research on LLMs 的四层 checklist}{Loubna Ben Allal 官方 deck，第 61 页}

\begin{importantbox}{开放 release 的最小定义}
\begin{enumerate}
  \item 能知道模型用了什么 data 和 transformation；
  \item 能获得或至少审计 weights/code/config；
  \item 能复现主要 training/evaluation path；
  \item 能看到 license、limitations、PII/opt-out 与 cost。
\end{enumerate}
仅有 downloadable weights，不足以回答这些问题。
\end{importantbox}

\teachervoice{00:43:10--00:45:10，讲者明确说 open release 不是“release model weights and stop there”，而是 data tools、opt-out、PII、reproducibility、evaluation tools 与 technical reports 的组合。}

\subsection{SantaCoder、StarCoder、StarCoder2：项目怎样迭代}

SantaCoder 作为约 1.1B ablation platform，StarCoder 扩到 15B/约 1T tokens，StarCoder2 提供 3B/7B/15B 并使用更大更多语言 corpus。slide 同时列出 Python source 比例下降，以及透明度与开放访问的延续。

\lecturefigure{slide-62.jpg}{从 SantaCoder 到 StarCoder2：规模、语言、Python 比例与开放性}{Loubna Ben Allal 官方 deck，第 62 页}

\begin{knowledgebox}{为什么先做 SantaCoder}
TB-scale data rule 无法直接在 15B full run 上反复试。小模型提供 filter/format/objective ablation platform；但 extrapolation 仍需在更大模型验证，不能假设所有 data effect 完全 scale-invariant。
\end{knowledgebox}

\subsection{Transparency Index：外部指标衡量披露，不衡量全部能力}

Stanford Foundation Model Transparency Index 按 data、labor、compute、risks、downstream 等披露维度评分。StarCoder 在课堂快照中排名靠前，说明文档/透明努力被外部框架捕捉；它不证明模型最准确、最安全或法律争议已解决。

\lecturefigure{slide-63.jpg}{Foundation Model Transparency Index 的课堂时结果}{Loubna Ben Allal 官方 deck，第 63 页}

\begin{warningbox}{Transparency 与 responsibility 相关但不等价}
披露坏结果比隐藏更透明，却仍可能不负责任；没有完整披露也可能有内部控制，但无法被公众验证。Index 衡量可见信息，不能替代 independent audit 或 legal judgment。
\end{warningbox}

\subsection{本章小结}

Responsible open model 需要 data、weights、code、evaluation、documentation 与 governance 一起发布。BigCode 的价值在于把这些层变成项目目标；透明度指标则提供外部检查，但不能替代能力与风险评估。

